\documentclass[11pt]{article}
\usepackage{amsmath,amssymb,amsthm,geometry,url}
\setlength{\emergencystretch}{2em}
\geometry{margin=1in}
\newtheorem{theorem}{Theorem}
\newtheorem{lemma}[theorem]{Lemma}
\newcommand{\Q}{\mathbb Q}\newcommand{\Z}{\mathbb Z}
\DeclareMathOperator{\Tr}{Tr}\DeclareMathOperator{\Nm}{Nm}
\title{A quartic in 198 variables}
\author{Research note; AI generated}
\date{14 November 2026}
\begin{document}\maketitle
\begin{abstract}
Under the arithmetic premises retained by Brox's 204-variable construction,
we obtain an integer quartic in 198 total variables, with 184 quadratic
squares and exactly 2945 monomials. An odd affine fourth label eliminates
its nonzero-index guard, and an inverse Vieta chart shortens the scalar
point block without an additional inverse witness. The enlarged point
domain admits the source unweighted trace-height bound. The sufficient
defining parameter remains ineffective. This is a six-variable reduction
from the pinned 204-variable baseline, not a large reduction in scale.
\end{abstract}

\section*{Statement and scope}
The source is Brox's 204-variable quartic, at repository commit
\begin{center}\small\ttfamily d337a5500e2531139445bd7e30f966f62c4db03a\end{center}
with precisely its
retained elliptic-group, scalar contact-height, pole obstruction and
prepared realization, modular contact-height, and fixed-data
simultaneous-prime premises. Its weighted elliptic trace-height theorem
is used only with constant positive weights. No additional $abc$
hypothesis is introduced and no independent verification of all these
arithmetic premises is asserted.

\begin{theorem}
Subject to those premises, the 198-coordinate quadratic table constructed
in \path{circuits/build_inverse198.py} has the following property. There is one fixed
integer $N_0$ such that, for every $N\ge N_0$, the rational input
projection of the common zero set of its residuals is exactly $\Z$.
The sufficient $N_0$ is ineffective. There are 198 total rational
coordinates and 184 quadratic residuals; summing their primitive integer
squares gives a polynomial of degree exactly four, content one, and
exactly 2945 nonzero coordinate monomials. Its coefficient height is
\[
 H(P_N)=\max\{C_0,149^{4N+2}\},\qquad
 C_0=(5898050446009651968\cdot31581660^{28})^2.
\]
The second term dominates for $N\ge53$. This is a coefficient crossover,
not a sufficient defining threshold. These exact algebraic assertions
are independently certified in
\path{circuits/inverse198/independent-verification.json};
the arithmetic projection assertion is proved below.
\end{theorem}

First exclude the zero stratum directly from the compiled rows, before
using any inverse rescaling. Let $J_s$ denote the rescaled scalar-product
coordinate (source coordinate 106 before compaction). The retained
input-square and scalar-product equations give
\[
 \begin{aligned}
RJ_s&=(51X^2+103)(51X^2+2503)\\
&\qquad\cdot(51X^2+2504)(1632X^2+80471)=r_s>0.
\end{aligned}
\]
Thus $R\ne0$. The retained forced-prime equation is
$1+22201RJ_s-149^{2N+1}J_s=0$, so
\[
 R=\frac{149^{2N+1}r_s}{1+22201r_s}>100 \qquad(N\ge1).
\]
Here each displayed factor is positive, and $r_s\ge1$ for every rational
input. The retained square row gives $R_2=R^2$. These facts justify all
subsequent Laurent inversions on every rational solution.

For clarity the proof uses physical coordinates, before the invertible
Laurent rescalings by the positive scalar normalizer. Those rescalings
and primitive integer row normalization are certified by the complete
table transformation. All variable counts include the input.

Put $K=\Q(\rho)$, $\rho^2=2$, $d=3-\rho$, $D=d^2$, and
$E:y^2=x^3-Dx$. Retain $P_0=2R_0$ and $2E(K)=\Z P_0$.
Every $Q\in E(K)$ is uniquely $[m]R_0+T_0$ with $T_0\in E[2](K)$;
write $\nu(Q)=m/2$. The source reference chart recovers $H,H-P_0$ and
$S=y(H)/x(H)$, $T_H=y(H-P_0)/x(H-P_0)$.
The scalar argument below excludes its forbidden torsion and reference
indices before any logarithm is used.

\section*{The inverse Vieta scalar chart}
Use field coordinates $U,V,W,M,r,P,r_2$, where $P$ here denotes a
coordinate, not the fixed point $P_0$. Write $R$ for the positive scalar
normalizer. The new equations are
\[
 UV=R^2P,\quad PW=r,\quad
 r(U+V+W+2M)+DR^2=0,\quad R^2P+W(U+V)=M^2,
\]
\[
 r_2=r^2,\qquad F_U r_2=U^2,\qquad F_V r_2=V^2,
\]
where $F_U,F_V$ reuse the old squared reciprocal-abscissa outputs. The
comparison and scalar equations are
\[
 L_2(M+U)=2RS,\qquad L_2^0(M+V)=-2RT_H,
 \qquad Z_{\rm st}(M+W)=-2S.
\]
The true scalar quotient is $Z=RZ_{\rm st}$.

Since $DR^2\ne0$, the third equation forces $r\ne0$. Consequently
$U,V,W,P$ are all nonzero. The ratio equations, once $S,T_H\ne0$ are
known, imply $M+U,M+V,M+W\ne0$. Define
\[
 U_2=\left(\frac{VW}{R^2},\frac{(M+U)VW}{R^3}\right),\qquad
 V_2=\left(\frac{UW}{R^2},-\frac{(M+V)UW}{R^3}\right),
\]
\[
                  Q_2=\left(P,-\frac{(M+W)P}{R}\right).
\]
The Vieta identities put all three points on $E$ and give
$Q_2=U_2-V_2$. Indeed the line $y=(M/R)x+r/R$ meets the cubic in
$U_2,-V_2,-Q_2$, with repeated intersections understood as tangencies.
The equations are precisely its three elementary symmetric identities.
All three ordinates are nonzero. Hence $Q_2$ is nontorsion, since all
$K$-rational torsion is two-torsion, and
\[
                         a=2\nu(Q_2)\in\Z\setminus\{0\}.
\]
The recovered reciprocal abscissas are $U/r$ and $V/r$, so the stored
$F_U,F_V$ are their actual squares. The displayed ratios give
\[
 L_2=2z(U_2)/z(H),\quad L_2^0=2z(V_2)/z(H-P_0),
 \qquad Z=2z(Q_2)/z(H),\quad z=-x/y.
\]

Conversely any finite nontorsion triple $U_2,V_2,Q_2=U_2-V_2$ supplies
this chart when its difference is nonzero. The secant through
$U_2,-V_2$ is nonvertical, including its tangent interpretation.
Its intercept is nonzero: a zero intercept would put $(0,0)$ among
its three intersection points, whereas all three are nontorsion.
Recovering its slope and intercept and applying the inverse scalings
therefore gives every displayed coordinate without a new guard.

\section*{The affine fourth label and the local factors}
Keep the source reciprocal pairs for $i=1,3$. For $i=4$ replace the
seven-field doubled-root block by the same five-field reciprocal pair
with coordinates $f_4,h_4,g_4,k_4,s_4$ and square outputs $F_4,G_4$:
\[
 F_4=f_4^2,\quad G_4=g_4^2,\quad h_4^2=f_4-Df_4F_4,
\]
\[
 s_4(f_4-g_4)=h_4+k_4,\qquad
 s_4(h_4-k_4)=1-D(F_4+f_4g_4+G_4).
\]
Its finite recovered difference $Q_4$ has
$x_4=s_4^2+D(f_4+g_4)$. Use the affine ratios
\[
 L_4=4z(U_4)/z(H)+1,\qquad
 L_4^0=4z(V_4)/z(H-P_0)+1,
\]
encoded by
\[
 h_4(L_4-1)+4f_4S=0,\quad(1-DF_4)(L_4-1)+4h_4S=0,
\]
\[
 k_4(L_4^0-1)+4g_4T_H=0,\quad
 (1-DG_4)(L_4^0-1)+4k_4T_H=0.
\]
These equations cover numerator identities, where the affine ratio is
$1$, and exclude nonidentity two-torsion numerator points exactly as in
the source. The fourth label is
\[
                       b=4\nu(Q_4)+1\in2\Z+1,
\]
so it is nonzero even if $Q_4$ is finite torsion. The other labels are
$n=\nu(Q_1)\in\tfrac12\Z$ and $t=2\nu(Q_3)\in\Z$.

All five local factors for each of the two detectors are retained, with
their reciprocal positivity tests and exponents. Their comparison
differences remain $C_i=L_i-\sigma L_i^0$. The input comparison is
$I=L_1-\beta$. Replace the Lucas expression by
\[
 G=L_3^2-(L_2+1)(L_2-1)L_3L_4+(L_2-1)^2L_4^2-1.
\]
Its two multiplication outputs are now $(L_2-1)L_4$ and
$(L_2+1)L_3$, so the affine change adds no coordinate.

\section*{Scalar extraction on the enlarged point domain}
The unchanged scalar rows recover
\[
 \beta=51X^2+2504,\quad r_s=\beta(\beta-1)(\beta-2401)(32\beta+343)>0,
\]
\[
 R=\frac{149^{2N+1}r_s}{1+149^2r_s}>100,
 \qquad F=\Nm(f(H))\ne0,
\]
and the exact physical quantities
\[
 W_s=\frac{Z_0^4}{R^{10}}+\frac{Z_1^2}{R^5}+\frac{F^2}{R^4},
 \qquad b_s=R(R+9792)W_s^2.
\]
In particular $b_s>0$. Let $p>76176$ or $p=149$, and suppose
$\mu=v_p(R)>0$ is odd. Every finite reciprocal-abscissa valuation is
even, so $v_p(F)$ is even. The three finite summand valuations in $W_s$
have respectively residues $2$ modulo $4$, odd residue modulo $2$,
and $0$ modulo $4$; zero terms are simply omitted. There is no
cancellation at their minimum. The scalar pole obstruction implies
$v_p(W_s)>-\mu/2$, and hence
\[
 v_p(Z_0)>19\mu/8,\qquad v_p(Z_1)>9\mu/4,
 \qquad v_p(F)>7\mu/4.
\]
The forced prime $149$ has the same positive odd normalizer valuation
$2N+1$ or $2N-1$ as in the source. The inequality for $F$ therefore
excludes reference torsion and the reference indices $-1,0,1$ using
exactly the source fixed-point norm computation. In particular
$S,T_H\ne0$, justifying all preceding point and ratio interpretations.

At some place above $p$, $f(H)$ has positive valuation, so $H$ is formal.
The new scalar quotient satisfies $v(Z)>9\mu/4$. Since $2$ is a unit,
$z(Q_2)$ is positive there. Thus $Q_2$ is formal or reduces to $(0,0)$;
a translation by that two-torsion point makes it formal and changes
$z$ only by sign. Write $\nu(H)=h$. Multiplying the group relation by
$2$ kills all torsion. The logarithmic ratio corresponding to $Z$ is
$2\nu(Q_2)/h=a/h$, up to this harmless sign. The formal logarithm is
an isometry, and $h\in\tfrac12\Z$ is $p$-integral. Therefore
\[
                         v_p(a)>9\mu/4.
\]
The source scalar radical and height argument applies without any
change of exponent or constant: the label is the integer $a$, and all
extra factors of $2$ are units at its consumed primes. With
$c=|a|$ and $L_N=\lfloor9(2N-1)/4\rfloor+1$ it gives
\[
 c\ge149^{L_N},\qquad
 h(\beta)\le C_3 149^{-L_N}c(1+\log c/3),\quad C_3=2^{330770}.
\]

\section*{The unweighted height aggregate and uniform soundness}
Use the quadratic aggregate
\[
 T=7+\frac14\Tr\bigl(x(Q_1)+P+x(Q_3)+x(Q_4)\bigr),
 \qquad t_0=\lambda,\quad t_1=\lambda+T.
\]
Every finite point has $x\ge-d$ at the corresponding real embedding,
so $T\ge1$. The source positive weighted trace-height theorem with all
weights $1$, together with $h(T-7)\le h(T)+\log8$, gives fixed constants
\[
 B^2\le A_0h(T)+J_0,\qquad
 B=\max(1,|n|,|a|,|t|,|b|).
\]
The factors $2$ and $4$ and the affine constant in the labels merely
change those fixed constants. All four points are finite, including any
torsion differences in the first, third, and fourth blocks.

At each consumed contact prime the unchanged local factors give the
same orders as before. Positive reciprocal tests put the numerator and
reference points in their required formal groups. The source transfers
integrality between the two full logarithmic ratios before applying
the relative order-four error as an absolute congruence. Multipliers
$2$ in block two and $4$ with an affine constant in block four are units
at these primes. Consequently
\[
 L_1\equiv n\pmod p,\quad L_2\equiv a\pmod{p^3},\quad
 L_3\equiv t\pmod{p^3},\quad L_4\equiv b\pmod{p^3}.
\]
Thus each contact radical $M_j=M(t_j)$ satisfies
\[
 M_j^3\mid\Delta,\qquad M_j\mid J_X\quad(J_X\ne0),
\]
where
\[
 \Delta=t^2-(a+1)(a-1)tb+(a-1)^2b^2-1\in\Z,
\]
and $J_X$ is the positive absolute reduced numerator of $\beta-n$.
Here $|\Delta|\le14B^4$ and $J_X\le4B e^{h(\beta)}$ when nonzero.
The two unchanged contact-height bounds and
$h(T)\le h(t_0)+h(t_1)+\log2$ imply: if $M_j^k\mid J\ne0$ for both
$j$, then
\[
                      B^2\le C|J|^{7/(6k)}+J_1.
\]
No weight-height term is needed in this construction.

If $\Delta\ne0$, the last inequality at $k=3$ gives
$B^2\le C_1B^{14/9}+J_1$, bounding $B$ by a fixed constant.
For all sufficiently large fixed $N$, this contradicts
$B\ge c\ge149^{L_N}$. Hence $\Delta=0$.
Put $A=a+1$, $c'=|A|$, and $w=\operatorname{sgn}(A)(A-2)b$.
For large $N$ these are well-defined with $w\ne0$, and
$t^2-c'tw+w^2=1$. The source norm-one-unit argument gives
$|w|=U_j$ and $|t|\in\{U_{j-1},U_{j+1}\}$ for the Lucas recurrence
with parameter $c'$. Divisibility by $|A-2|$ gives $j\ge c'-2$.
In particular $B\ge|t|\ge U_{j-1}\ge(c'-1)^{j-2}$, so
\[
             \log B\ge(c'-5)\log(c'-2)
                     \ge(c-6)\log(c-3)
\]
for sufficiently large $c$.

If $J_X\ne0$, the transfer bound with $k=1$, once $B$ exceeds a fixed
constant, gives $\log B\le(7/5)h(\beta)+C_2$.
Choose $N$ large enough that $(7/5)C_3 149^{-L_N}\le1$.
Its right side is at most $c+(c/3)\log c+C_2$, contradicting the
preceding lower bound for sufficiently large $c$. Thus
$\beta=n\in\tfrac12\Z$. For reduced $X=u/v$, this implies
$v^2\mid102$, and squarefreeness gives $v=1$. All thresholds are
fixed independently of the input and witnesses.

\section*{Uniform scalar completeness after multiplying the quotient by two}
Fix integer $X$. The unchanged four-factor normalizer has $r_s\ge1$,
so $R\ge149^{2N+1}/(1+149^2)$ uniformly in $X$.
Use the fixed real neighborhoods of source P13, where the old quotient
$Z^{\rm old}=z(Q_2)/z(H)$ and $F$ satisfy
$|Z^{\rm old}_0|,|Z^{\rm old}_1|\le M_Z$ and $|F|\le M_F$.
Since the new quotient is $2Z^{\rm old}$, for $R\ge1$ one has
\[
 0<W_s\le K'_W R^{-4},\qquad
 K'_W=16M_Z^4+4M_Z^2+M_F^2,
\]
\[
                   0<b_s\le9793(K'_W)^2 R^{-6}.
\]
Increasing a fixed scalar threshold therefore gives $b_s<17$ for every
integer input and all larger $N$.

For a particular $(N,X)$, follow the source P14 finite-place preparation:
choose $H$ sufficiently formal at the finite support and then an
arbitrarily large positive integer $a_0$, sufficiently divisible and
with $[a_0]P_0$ in the fixed real neighborhood of $O$. Impose in addition
$3\mid a_0$. Compact-group recurrence in finite-index progressions
permits this extra fixed divisibility. Set
\[
                     Q_2=[a_0]P_0=[a]R_0,\qquad a=2a_0\in6\Z.
\]
The new factor $2$ in $Z$ has zero valuation at every odd prime and
increases its $2$-adic component depths. The three individual summands
of $W_s$ can still be made integral at its odd support and arbitrarily
deep at $2,3,17$ by the same ordered choices of $H$ and $a_0$.
The squareclass identity $b_s=R(R+9792)W_s^2$ is unchanged, as are the
normalizer's square denominator and its local squareclasses. Thus the
prepared pole-realization premise applies to the actual new scalar input.
After $a_0$ is fixed, continuity preserves these conditions for $H$ in
the source progression $h\equiv h_*\pmod L$, with $64\mid h_*,L$.
No input-dependent lower bound on $N$ has been introduced.

\section*{Completion of the fresh witnesses}
For the chosen $a\in6\Z$, set $A=a+1$, $j=a-1$, and
\[
                   b=U_j(A)/(a-1),\qquad t=U_{j+1}(A).
\]
Modulo $a-1=A-2$, $U_j\equiv j$, so $b$ is integral.
Modulo $2$ the recurrence for odd $A$ has period $0,1,1$;
$j\equiv2\pmod3$ and $a-1$ is odd, so $b$ is odd. Choosing $a$
large makes $b>1$. The Lucas identity gives $\Delta=0$.
Take $Q_1=[\beta]P_0$, $Q_3=[t]R_0$ and
$Q_4=[(b-1)/2]R_0$. These points, $Q_2$, and $T$ are now fixed
before comparison primes are chosen.

For integer reference $h$ use $H=[h]P_0$ and
$U_i=[h]Q_i$, $V_i=[h-1]Q_i$. All scalar-chart numerator points
are nontorsion and nonzero when $h\ne0,1$, and their difference is
$Q_2\ne O$; hence the inverse Vieta converse applies.
The source Gaussian reciprocity lemma annihilates $R_0$ at the same
conjugate order moduli as before, so the same compatible $h\equiv0$
and $h\equiv1$ congruences put every new numerator pair in its required
formal groups. The scaled and affine ratios have exact limiting labels
$\beta,a,t,b$. The same compact-group recurrence supplies the real
neighborhoods, and the same fixed-data simultaneous-prime premise
applies to $\Theta\cup(\Theta-T)$, with repeated coefficients merged
and incidences retained.

The unchanged contact equations and factor exponents then give every
nonconstant factor summand positive valuation at a selected place, and
reciprocal-trace residue $1$ or $2$ as in the source. At prepared primes
the weights tend to zero. On fixed real neighborhoods the weights tend
to zero positively, and each comparison pole input tends to
$(408/409)^2/4\in(0,17)$. Its reduced denominator is square and its
local squareclasses and depths are the source ones. Apply the prepared
pole realization to these actual inputs and fill all chart coordinates.
The invertible rescalings give a common zero of all 184 residuals.
This proves completeness for every integer input beyond one fixed scalar
threshold. Combining it with soundness proves the theorem.

\section*{Balanced coefficients and reproducibility}
There is also an invertibly rescaled presentation with the same 198
variables, 184 quadratic residuals and 2945 coordinate monomials.
Put $C=149^N$ and $K_0=31581660^{28}$. Its exact coefficient height is
\[
 H_{\mathrm{bal}}(P_N)=\max\{C_0,K_0\,149^{2N}\}.
\]
The second term dominates for $N\ge57$; this again is only a coefficient
threshold. Every rescaling multiplier is a nonzero power of the fixed
integer $C$, so the two presentations have the same rational input
projection at every fixed $N$. This preserves the coefficient-height
bound of the 204-variable baseline while reducing its support from
3046 to 2945 monomials.

The complete equations and expanded polynomial are the files
\path{circuits/inverse198/balanced-quadratic.json} and
\path{circuits/inverse198/balanced-polynomial.json}.
The producer is \path{circuits/balance_inverse198.py}; the
independent checker \path{circuits/check_balanced198.py} verifies all
184 row rescalings and expands their squares separately. It checks the
coefficient maximum and proves that the support does not change at any
$C=149^N$, $N\ge1$. The unbalanced system has two further independent
reconstruction checkers. These exact computations concern algebraic
identities and coefficient assertions; they do not certify the external
arithmetic premises.

\section*{Dependencies and finite verification}
The source scalar assertions used above are P04 and P11--P14 of the
pinned upstream \path{research/quartic214/proofs/uniform-existence.md},
commit
\begin{center}\small\ttfamily 1f7e1f1fb0438feda3c9f3da9cc16982126aeeed\end{center}
The source contact, formal-group, and prime applications are those of
Brox's 204-variable proof, with its separately proved weighted trace
and half-generator lemmas. The new complete chart identities, all
Laurent row transformations, and the full quartic expansion require their
separate exact certificates. The proof does not identify any small
numerical value of $N$ as defining the integers.
\section*{Larger reductions examined}
The reduction above is modest. Several attempts to remove an entire
point or detector block did not produce a valid improvement. The
following obstructions are retained to distinguish those attempts from
the proved construction.

\textbf{Eliminating the input point.}
Replacing the Pell coordinate by $w=(a-2)b+\beta$ makes
$t^2-atw+w^2-1$ monic in $\beta$. If this expression vanished with
integer $a,t,b$, rational-root integrality would give $\beta\in\Z$
and remove the separate input-point comparison. Before vanishing,
however, clearing the denominator introduces a factor of order
$\exp(2h(\beta))$. The current height bounds yield only
$\log B=O(h(\beta))$, which is compatible with the scalar estimate.
The exponential Lucas lower bound becomes available only after exact
vanishing. This leaves a substantive gap, not a completed reduction.

\textbf{Merging the detectors.}
The contact function satisfies $R(2-t)=R(t)$. On the admissible reflection
locus $\lambda=1-T/2$, the two arguments are reflections, so
$R(\lambda)R(\lambda+T)=R(\lambda)^2$. Multiplication erases all odd
valuations. Both contact values are positive for sufficiently large
$T$, so real positivity does not repair this loss.

\textbf{Prescribing the scalar index.}
A polynomial such as $149^{kN}r_s^3$ supplies large prime valuations,
but it need not be divisible by the elliptic reduction orders required
to make the scalar point formal at those primes. Thus it cannot simply
replace the source's freely chosen sufficiently divisible index.

\textbf{Other exact obstructions.}
The accompanying research notes give prepared-input counterexamples to
two proposed smaller pole predicates. They also record the unwanted
local four-torsion branch in the parameter $z([2]Q)/2$. A full-rank
finite-field certificate rules out a quadratic relation among the four
unchanged reference outputs $S,T_H,f(H)^2,f(H-P_0)^2$; this obstructs
one direct coordinate deletion, not every possible reference chart.
No global optimality claim is made.
\end{document}
